% Source: ACM_Templates.pdf, PDF p. 81.
\subsection{树的重心}
\label{tree:centroid}
\AlgorithmInfo
    {找删除后最大连通块点数最小的顶点，用于平衡划分、点分治等。}
    {时间、空间 $O(n)$；递归深度最坏 $O(n)$。}
    {按顶点个数计算，非空连通树；无向边只添加一次。}
\subsubsection*{原理}
任意选根，删除 $u$ 后的连通块由各儿子子树，以及大小为 $n-sz_u$ 的父亲方向部分构成。因此
\[
    balance_u=\max\left(n-sz_u,\ \max_{v\text{ 为 }u\text{ 的儿子}}sz_v\right).
\]
重心等价于 $balance_u\le n/2$。若有大于一半的连通块，向它移动会改善最坏大小；不存在这样的块时已经平衡。树至少有一个重心，至多两个，两个重心一定相邻；这里是等点权计数，不是带点权重心。
\subsubsection*{使用教程}
\InlineCode{init(n)} 后用 \InlineCode{add(u,v)} 加无向边，每边一次；\InlineCode{build(root)} 计算结果，省略根取 1。\InlineCode{centroid} 是全部重心，\InlineCode{balance[u]} 为删 $u$ 后的最大块大小，\InlineCode{sz[u]} 为当前根下的子树大小。重跑保留边、重置结果；选根不改变重心集合。
\begin{lstlisting}[style=mtycpp,numbers=none]
TreeCentroid g;
g.init(4);
g.add(1,2),g.add(2,3),g.add(3,4);
g.build(1);
auto centroid=g.centroid;
int largest=g.balance[2];
\end{lstlisting}
\needspace{7\baselineskip}
\subsubsection*{模板代码}
\CodeFile{code/trees/centroid.cpp}
